% =========================================================
% ANEXOS
% FINALIDADE: reunir material NÃO ELABORADO PELO(A) AUTOR(A), incluído
%   para fundamentação, comprovação ou ilustração (por exemplo, o parecer
%   do Comitê de Ética em Pesquisa). Material próprio vai em apendices.tex.
% ESCREVA: cada anexo começa com \anexo{TÍTULO EM MAIÚSCULAS}.
%   A letra (A, B, C...) e a entrada no sumário são automáticas.
% OPCIONAL: sem anexos, apague este arquivo; o elemento é omitido.
% CUIDADO: cite cada anexo no texto, por exemplo:
%   conforme o Anexo~\ref{anx:exemplo}.
% EXEMPLO de segundo anexo (retire os % para usar):
%   \anexo{PARECER DO COMITÊ DE ÉTICA EM PESQUISA}
%   \label{anx:parecer}
%   Texto do anexo.
% =========================================================

\anexo{TÍTULO DO ANEXO}
\label{anx:exemplo}

Elemento opcional. Documento não elaborado pelo(a) autor(a), que serve de fundamentação, comprovação ou ilustração.

% Para inserir um documento em PDF, salve-o na pasta 5-figuras/ e
% retire o % da linha abaixo, ajustando o nome do arquivo.
% pages=- insere todas as páginas; pages={1-3} insere só as páginas 1 a 3.
% \includepdf[pages=-,pagecommand={}]{5-figuras/nome-do-arquivo.pdf}
